\section{Introduction}
\label{sec:intro}

Current and upcoming CMB experiments (Simons Observatory, CMB-S4, CCAT) and
large-scale-structure surveys interpret their data through cross-correlations
of lensing, the thermal and kinetic Sunyaev--Zel'dovich (tSZ/kSZ) effects, and
the cosmic infrared background (CIB) with galaxy and cluster catalogues. These
analyses need many-realization mock skies: correlated maps of all secondary
anisotropies together with the halo catalogues that source them, for pipeline
validation, covariance estimation, and foreground-model exploration. Full
$N$-body lightcones with gas physics remain far too expensive for ensembles of
hundreds of realizations, which motivates approximate methods calibrated
against them.

The mass-peak-patch algorithm \citep{BondMyers1996a, Stein2019} occupies a
productive point on the accuracy--cost curve: haloes are identified as peaks of
the smoothed initial density field using homogeneous-ellipsoid collapse,
exclusion-and-merging resolves overlapping proto-haloes, and second-order
Lagrangian perturbation theory (2LPT) moves them to their final positions and
velocities. Its best-known product is the Websky simulation
\citep{Stein2020}, whose full-sky halo catalogue and correlated
$\kappa$/kSZ/tSZ/CIB maps have become a community standard for
secondary-anisotropy work. Related mock-sky efforts include Sehgal et
al.-style template skies \citep{Sehgal2010}, the \pkg{Agora} multi-component
sky \citep{Omori2024}, and the HalfDome $N$-body lightcones
\citep{HalfDome2024}; all such suites, like peak patch, calibrate their halo
masses by abundance matching to a target mass function.

Two practical constraints of the original Fortran implementation motivate this
work. First, the peak-patch computation is dominated by large distributed FFTs
of the initial conditions and by the per-filter collapse searches; on CPU
clusters a Websky-resolution run occupies hundreds to thousands of cores, with
the internode FFT transpose as the scaling bottleneck. Second, ensemble
applications --- for example the planned $\gtrsim$1000-realization Backlight
ensemble at CITA --- put a premium on time-to-solution \emph{per node}, where
modern GPU nodes offer an order of magnitude more arithmetic throughput than
the CPU nodes the original code targets.

\ppjl\ is a from-scratch Julia implementation of the algorithm designed around
these constraints. Its central algorithmic ingredient
(Section~\ref{sec:multires}) is a multi-resolution decomposition of the
initial Gaussian random field, in the spirit of the multi-scale initial
conditions of \pkg{MUSIC} \citep{HahnAbel2011}: a global coarse grid carries
the long-wavelength modes, and each spatial tile synthesizes its own fine
modes from a counter-based random number generator keyed to the global lattice
site, so that the union of tiles reproduces a single global realization
exactly, with no global fine-grid FFT and no internode communication. Each
tile then runs the full peak-patch pipeline independently on a GPU
(Section~\ref{sec:implementation}), which makes the cost embarrassingly
parallel over tiles and lets a single node process a $6144^3$-equivalent
volume.

Beyond the halo-finding pipeline, \ppjl\ integrates map making
(Section~\ref{sec:maps}): line-of-sight ``field'' kernels painted directly
from the Lagrangian lattice (lensing convergence, electron-density-weighted
momentum for kSZ, and the linear integrated Sachs--Wolfe effect), and
halo-profile painting (NFW convergence, Battaglia-type gas density and
pressure, CIB halo model) through the \pkg{XGPaint} package
\citep{XGPaint}. The combination reproduces the construction of the released
Websky maps, which we use as our end-to-end validation target
(Sections~\ref{sec:valcat}--\ref{sec:valmaps}): six catalogue-level statistics
match at the few-per-cent level, and the map-level angular power spectra match
the released products over their well-measured multipole ranges. The
validation campaign also surfaced a resolution criterion specific to
velocity-weighted maps --- the coarse grid must resolve the velocity field to
$k_{\rm Nyq}\gtrsim 0.3\,h\,{\rm Mpc}^{-1}$ --- which we derive and verify in
Section~\ref{sec:multires-velocity}.

This paper is organized as follows. Section~\ref{sec:algorithm} reviews the
peak-patch algorithm and fixes conventions. Section~\ref{sec:multires}
describes the multi-resolution field synthesis. Section~\ref{sec:implementation}
details the GPU implementation and the MPI alternative.
Section~\ref{sec:lightcone} describes lightcone construction and catalogue
products. Section~\ref{sec:maps} presents the map-making machinery.
Sections~\ref{sec:valcat} and \ref{sec:valmaps} validate catalogues and maps
against Websky. Section~\ref{sec:performance} reports performance, and
Section~\ref{sec:summary} concludes.
